\section{Finding suitable affine opens}
\label{section-finding-affine-opens}

\noindent
In this section we collect some results on the existence of
affine opens in more and less general situations.

\begin{lemma}
\label{lemma-maximal-points-affine}
Let $X$ be a quasi-separated scheme.
Let $Z_1, \ldots, Z_n$ be pairwise distinct irreducible components of $X$,
see Topology, Section \ref{topology-section-irreducible-components}.
Let $\eta_i \in Z_i$ be their generic points, see
Schemes, Lemma \ref{schemes-lemma-scheme-sober}.
There exist affine open neighbourhoods $\eta_i \in U_i$
such that $U_i \cap U_j = \emptyset$ for all $i \not = j$.
In particular, $U = U_1 \cup \ldots \cup U_n$ is an affine
open containing all of the points $\eta_1, \ldots, \eta_n$.
\end{lemma}

\begin{proof}
Let $V_i$ be any affine open containing $\eta_i$
and disjoint from the closed set $Z_1 \cup \ldots \hat Z_i \ldots \cup Z_n$.
Since $X$ is quasi-separated for each $i$ the union
$W_i = \bigcup_{j, j \not = i} V_i \cap V_j$ is a quasi-compact
open of $V_i$ not containing $\eta_i$.
We can find open neighbourhoods $U_i \subset V_i$
containing $\eta_i$ and disjoint from $W_i$ by
Algebra, Lemma \ref{algebra-lemma-standard-open-containing-maximal-point}.
Finally, $U$ is affine since it is the spectrum of
the ring $R_1 \times \ldots \times R_n$ where $R_i = \mathcal{O}_X(U_i)$,
see Schemes, Lemma \ref{schemes-lemma-disjoint-union-affines}.
\end{proof}

\begin{remark}
\label{remark-maximal-points-affine}
Lemma \ref{lemma-maximal-points-affine} above is false if $X$
is not quasi-separated. Here is an example. Take
$R = \mathbf{Q}[x, y_1, y_2, \ldots]/((x-i)y_i)$.
Consider the minimal prime ideal $\mathfrak p = (y_1, y_2, \ldots)$
of $R$. Glue two copies of $\Spec(R)$ along the
(not quasi-compact) open $\Spec(R) \setminus V(\mathfrak p)$
to get a scheme $X$ (glueing as in
Schemes, Example \ref{schemes-example-affine-space-zero-doubled}).
Then the two maximal points of $X$ corresponding to $\mathfrak p$
are not contained in a common affine open. The reason is
that any open of $\Spec(R)$ containing $\mathfrak p$
contains infinitely many of the ``lines'' $x = i$, $y_j = 0$,
$j \not = i$ with parameter $y_i$. Details omitted.
\end{remark}

\noindent
Notwithstanding the example above, for ``most'' finite sets of irreducible
closed subsets one can apply Lemma \ref{lemma-maximal-points-affine} above,
at least if $X$ is quasi-compact. This is true because $X$ contains a dense
open which is separated.

\begin{lemma}
\label{lemma-quasi-compact-dense-open-separated}
Let $X$ be a quasi-compact scheme.
There exists a dense open $V \subset X$ which is separated.
\end{lemma}

\begin{proof}
Say $X = \bigcup_{i = 1, \ldots, n} U_i$ is a union of $n$ affine open
subschemes. We will prove the lemma by induction on $n$. It is trivial for
$n = 1$. Let $V' \subset \bigcup_{i = 1, \ldots, n - 1} U_i$ be a separated
dense open subscheme, which exists by induction hypothesis. Consider
$$
V = V' \amalg (U_n \setminus \overline{V'}).
$$
It is clear that $V$ is separated and a dense open subscheme of $X$.
\end{proof}

\noindent
It turns out that, even if $X$ is quasi-separated as well as quasi-compact,
there does not exist a separated, quasi-compact dense open, see
Examples, Lemma
\ref{examples-lemma-no-dense-separated-quasi-compact-open-in-qcqs}.
Here is a slight refinement of Lemma \ref{lemma-maximal-points-affine} above.

\begin{lemma}
\label{lemma-point-and-maximal-points-affine}
Let $X$ be a quasi-separated scheme. Let $Z_1, \ldots, Z_n$ be pairwise
distinct irreducible components of $X$. Let $\eta_i \in Z_i$ be their
generic points. Let $x \in X$ be arbitrary.
There exists an affine open $U \subset X$ containing
$x$ and all the $\eta_i$.
\end{lemma}

\begin{proof}
Suppose that $x \in Z_1 \cap \ldots \cap Z_r$ and
$x \not \in Z_{r + 1}, \ldots, Z_n$. Then we may choose
an affine open $W \subset X$ such that
$x \in W$ and $W \cap Z_i = \emptyset$ for
$i = r + 1, \ldots, n$. Note that clearly $\eta_i \in W$
for $i = 1, \ldots, r$. By Lemma \ref{lemma-maximal-points-affine}
we may choose affine opens $U_i \subset X$ which are pairwise
disjoint such that $\eta_i \in U_i$ for $i = r + 1, \ldots, n$.
Since $X$ is quasi-separated the opens $W \cap U_i$
are quasi-compact and do not contain $\eta_i$ for
$i = r + 1, \ldots, n$. Hence by
Algebra, Lemma \ref{algebra-lemma-standard-open-containing-maximal-point}
we may shrink $U_i$ such that $W \cap U_i = \emptyset$
for $i = r + 1, \ldots, n$. Then the union
$U = W \cup \bigcup_{i = r + 1, \ldots, n} U_i$ is disjoint and hence
(by Schemes, Lemma \ref{schemes-lemma-disjoint-union-affines})
a suitable affine open.
\end{proof}

\begin{lemma}
\label{lemma-ample-finite-set-in-affine}
Let $X$ be a scheme. Assume either
\begin{enumerate}
\item The scheme $X$ is quasi-affine.
\item The scheme $X$ is isomorphic to a locally closed subscheme
of an affine scheme.
\item There exists an ample invertible sheaf on $X$.
\item The scheme $X$ is isomorphic to a locally closed subscheme
of $\text{Proj}(S)$ for some graded ring $S$.
\end{enumerate}
Then for any finite subset $E \subset X$ there exists an
affine open $U \subset X$ with $E \subset U$.
\end{lemma}

\begin{proof}
By Properties, Definition \ref{definition-quasi-affine}
a quasi-affine scheme is a quasi-compact open subscheme
of an affine scheme. Any affine scheme $\Spec(R)$ is isomorphic to
$\text{Proj}(R[X])$ where $R[X]$ is graded by setting $\deg(X) = 1$.
By Proposition \ref{proposition-characterize-ample}
if $X$ has an ample invertible sheaf then $X$ is isomorphic to an open
subscheme of $\text{Proj}(S)$ for some graded ring $S$.
Hence, it suffices to prove the lemma in case (4).
(We urge the reader to prove case (2) directly for themselves.)

\medskip\noindent
Thus assume $X \subset \text{Proj}(S)$ is a locally closed subscheme
where $S$ is some graded ring. Let $T = \overline{X} \setminus X$.
Recall that the standard opens $D_{+}(f)$ form a basis of the
topology on $\text{Proj}(S)$. Since $E$ is finite we may choose finitely many
homogeneous elements $f_i \in S_{+}$ such that
$$
E \subset
D_{+}(f_1) \cup \ldots \cup D_{+}(f_n) \subset
\text{Proj}(S) \setminus T
$$
Suppose that $E = \{\mathfrak p_1, \ldots, \mathfrak p_m\}$
as a subset of $\text{Proj}(S)$.
Consider the ideal $I = (f_1, \ldots, f_n) \subset S$.
Since $I \not \subset \mathfrak p_j$ for all $j = 1, \ldots, m$
we see from Algebra, Lemma \ref{algebra-lemma-graded-silly} that
there exists a homogeneous element $f \in I$, $f \not \in \mathfrak p_j$
for all $j = 1, \ldots, m$. Then $E \subset D_{+}(f) \subset
D_{+}(f_1) \cup \ldots \cup D_{+}(f_n)$. Since $D_{+}(f)$ does not
meet $T$ we see that $X \cap D_{+}(f)$ is a closed subscheme of the
affine scheme $D_{+}(f)$, hence is an affine open of $X$ as desired.
\end{proof}

\begin{lemma}
\label{lemma-ample-finite-set-in-principal-affine}
Let $X$ be a scheme. Let $\mathcal{L}$ be an ample invertible sheaf on $X$.
Let
$$
E \subset W \subset X
$$
with $E$ finite and $W$ open in $X$. Then there exists an $n > 0$
and a section $s \in \Gamma(X, \mathcal{L}^{\otimes n})$ such that
$X_s$ is affine and $E \subset X_s \subset W$.
\end{lemma}

\begin{proof}
The reader can modify the proof of Lemma \ref{lemma-ample-finite-set-in-affine}
to prove this lemma; we will instead deduce the lemma from it.
By Lemma \ref{lemma-ample-finite-set-in-affine} we can choose an affine
open $U \subset W$ such that $E \subset U$.
Consider the graded ring $S = \Gamma_*(X, \mathcal{L}) =
\bigoplus_{n \geq 0} \Gamma(X, \mathcal{L}^{\otimes n})$.
For each $x \in E$ let $\mathfrak p_x \subset S$ be the graded ideal
of sections vanishing at $x$. It is clear that $\mathfrak p_x$ is
a prime ideal and since some power of $\mathcal{L}$ is globally
generated, it is clear that $S_{+} \not \subset \mathfrak p_x$.
Let $I \subset S$ be the graded ideal of sections vanishing on all
points of $X \setminus U$. Since the sets $X_s$ form a basis
for the topology we see that $I \not \subset \mathfrak p_x$ for
all $x \in E$.
By (graded) prime avoidance (Algebra, Lemma \ref{algebra-lemma-graded-silly})
we can find $s \in I$ homogeneous
with $s \not \in \mathfrak p_x$ for all $x \in E$.
Then $E \subset X_s \subset U$ and $X_s$ is affine by
Lemma \ref{lemma-affine-cap-s-open}.
\end{proof}

\begin{lemma}
\label{lemma-quasi-affine-invertible-nonvanishing-section}
Let $X$ be a quasi-affine scheme. Let $\mathcal{L}$ be an invertible
$\mathcal{O}_X$-module. Let $E \subset W \subset X$ with $E$ finite
and $W$ open. Then there exists an $s \in \Gamma(X, \mathcal{L})$
such that $X_s$ is affine and $E \subset X_s \subset W$.
\end{lemma}

\begin{proof}
The proof of this lemma has a lot in common with the proof of
Algebra, Lemma \ref{algebra-lemma-silly}.
Say $E = \{x_1, \ldots, x_n\}$. If $E = W = \emptyset$, then $s = 0$
works. If $W \not = \emptyset$, then we may assume $E \not = \emptyset$
by adding a point if necessary. Thus we may assume $n \geq 1$.
We will prove the lemma by induction on $n$.

\medskip\noindent
Base case: $n = 1$. After replacing $W$ by an affine open neighbourhood
of $x_1$ in $W$, we may assume $W$ is affine. Combining
Lemmas \ref{lemma-quasi-affine-O-ample} and
Proposition \ref{proposition-characterize-ample}
we see that every quasi-coherent
$\mathcal{O}_X$-module is globally generated.
Hence there exists
a global section $s$ of $\mathcal{L}$ which does not vanish at $x_1$.
On the other hand, let $Z \subset X$ be the
reduced induced closed subscheme on $X \setminus W$.
Applying global generation to the quasi-coherent ideal sheaf
$\mathcal{I}$ of $Z$ we find a global section $f$ of $\mathcal{I}$
which does not vanish at $x_1$. Then $s' = fs$ is a global section
of $\mathcal{L}$ which does not vanish at $x_1$ such that
$X_{s'} \subset W$. Then $X_{s'}$ is affine by
Lemma \ref{lemma-affine-cap-s-open}.

\medskip\noindent
Induction step for $n > 1$. If there is a specialization
$x_i \leadsto x_j$ for $i \not = j$, then it suffices
to prove the lemma for $\{x_1, \ldots, x_n\} \setminus \{x_i\}$
and we are done by induction. Thus we may assume there are no
specializations among the $x_i$.
By either Lemma \ref{lemma-ample-finite-set-in-affine} or
Lemma \ref{lemma-ample-finite-set-in-principal-affine}
we may assume $W$ is affine.
By induction we can find a global section
$s$ of $\mathcal{L}$ such that $X_s \subset W$ is affine and contains
$x_1, \ldots, x_{n - 1}$. If $x_n \in X_s$ then we are done.
Assume $s$ is zero at $x_n$. By the case $n = 1$ we can find
a global section $s'$ of $\mathcal{L}$ with
$\{x_n\} \subset X_{s'} \subset
W \setminus \overline{\{x_1, \ldots, x_{n - 1}\}}$.
Here we use that $x_n$ is not a specialization of $x_1, \ldots, x_{n - 1}$.
Then $s + s'$
is a global section of $\mathcal{L}$ which is nonvanishing
at $x_1, \ldots, x_n$ with $X_{s + s'} \subset W$ and
we conclude as before.
\end{proof}

\begin{lemma}
\label{lemma-ring-affine-open-injective-local-ring}
Let $X$ be a scheme and $x \in X$ a point. There exists an affine open
neighbourhood $U \subset X$ of $x$ such that the canonical map
$\mathcal{O}_X(U) \to \mathcal{O}_{X, x}$ is injective in each of
the following cases:
\begin{enumerate}
\item $X$ is integral,
\item $X$ is locally Noetherian,
\item $X$ is reduced and has a finite number of irreducible components.
\end{enumerate}
\end{lemma}

\begin{proof}
After translation into algebra, this follows from
Algebra, Lemma \ref{algebra-lemma-subring-of-local-ring}.
\end{proof}

\begin{lemma}
\label{lemma-standard-open-two-affines}
Let $U$, $V$ be affine schemes and let $W \to U$ and $W \to V$
be open immersions. For any $w \in W$ there exists an affine
open neighbourhood $W' \subset W$ of $w$ such that $W'$
maps to a standard open of both $U$ and $V$.
\end{lemma}

\begin{proof}
Let $X$ be the glueing of $U$ and $V$ along $W$ and
apply Schemes, Lemma \ref{schemes-lemma-standard-open-two-affines}.
\end{proof}






